\documentclass[11pt]{article}
\usepackage[T1]{fontenc}
\usepackage[utf8]{inputenc}
\usepackage{lmodern,amsmath,amssymb,amsthm,mathtools}
\usepackage[margin=28mm]{geometry}
\usepackage{microtype,xcolor}
\usepackage[authoryear,round]{natbib}
\usepackage{xurl}
\usepackage[colorlinks=true,allcolors=blue!50!black]{hyperref}
\newtheorem{theorem}{Theorem}[section]
\newtheorem{proposition}[theorem]{Proposition}
\newtheorem{lemma}[theorem]{Lemma}
\newtheorem{corollary}[theorem]{Corollary}
\theoremstyle{definition}\newtheorem{definition}[theorem]{Definition}
\theoremstyle{remark}\newtheorem{remark}[theorem]{Remark}
\newcommand{\OO}{\mathcal O}
\newcommand{\ZZ}{\mathbf Z}
\newcommand{\QQ}{\mathbf Q}
\newcommand{\FF}{\mathbf F}
\newcommand{\dR}{\mathrm{dR}}
\newcommand{\cris}{\mathrm{cris}}
\newcommand{\et}{\mathrm{et}}
\newcommand{\Tr}{\operatorname{Tr}}
\newcommand{\Lie}{\operatorname{Lie}}
\newcommand{\End}{\operatorname{End}}
\newcommand{\Hom}{\operatorname{Hom}}
\newcommand{\Pic}{\operatorname{Pic}}
\newcommand{\Spf}{\operatorname{Spf}}
\newcommand{\Spec}{\operatorname{Spec}}
\newcommand{\workbenchurl}{https://github.com/KokunoYumeto/schemes-and-moduli-workbench}
\newcommand{\recordurl}{https://github.com/KokunoYumeto/mathematics-workbenches}
\title{Grothendieck's Hasse and ordinary determinant duality, after Bijakowski and Serre--Tate}
\author{}
\date{23 November 2026}
\begin{document}
\maketitle
\begin{abstract}
Grothendieck's crystalline notes compare the Hodge determinant of an abelian scheme with that of its dual, and ask how this comparison acts on the Hasse invariant and on the ordinary \'etale Tate determinant. We identify the trace normalization of the known Hasse-duality theorem of Bijakowski over every characteristic-$p$ base. On the ordinary locus, the determinant of a polarization divided by the positive square root of its degree is an integral isomorphism, independent of the polarization. With positive Poincar\'e biduality it becomes its inverse on the dual abelian scheme. Its reduction agrees with the inverse Hodge--trace comparison up to $(-1)^{g(g+1)/2}$. The product rule contains the sign $(-1)^{gh}$. The arguments include characteristic two, nonreduced bases, and polarizations whose degrees are divisible by $p$. The proof uses Serre--Tate canonical lifting, Riemann--Roch, and an explicit cup-product calculation. The mathematics and text were produced by GPT-6 Astra Pro and checked by the same model through complete proof and source comparison, separate statement-only reconstructions, and exact calculations; no human has checked them.
\end{abstract}

\section{Grothendieck's questions and the two comparisons}\label{sec:intro}

This paper makes two duality comparisons in Grothendieck's crystalline notes precise enough to use integrally and in families. On archive sheet~21 of cote~9, he identifies the Hodge determinant lines by the top de Rham trace and says that the isomorphism ``donc transforme $\sigma_A$ en $\sigma_{A^*}$'' \citep[cote~9, sheet~21, PDF image~22]{Grothendieck}. On sheet~28 he writes the inverse-duality formula $\alpha_{A^*}=\alpha_A^{-1}$, notes a sign question, and marks the ordinary de Rham comparison ``\`a v\'erifier'' \citep[cote~9, sheet~28, PDF image~29]{Grothendieck}. These are questions about canonical lines, rather than just equality of the ordinary loci of an abelian scheme and its dual. A scalar or a sign left undetermined in such a comparison changes its compatibility with biduality and products.

We fix the conventions first. Let $p$ be any prime, $S$ any $\FF_p$-scheme, and $f:A\to S$ an abelian scheme of constant relative dimension $g\ge1$. No reducedness, noetherian, or perfection hypothesis is imposed on $S$. Write $A^*$ for the dual abelian scheme and use the positive Poincar\'e biduality $\iota_A:A\xrightarrow{\sim}A^{**}$. For an ample line bundle $L$, our polarization is
\begin{equation}\label{eq:polarization-convention}
 \lambda_L(x)=t_x^*L\otimes L^{-1}.
\end{equation}
Thus $\lambda_L^\vee\iota_A=\lambda_L$. Set
\[
 W_A=f_*\Omega^1_{A/S},\qquad \omega_A=\det W_A,
 \qquad H_A=H^1_{\dR}(A/S).
\]
The Hodge exact sequence is
\begin{equation}\label{eq:hodge}
 0\longrightarrow W_A\longrightarrow H_A
 \longrightarrow W_{A^*}^{\vee}\longrightarrow0.
\end{equation}
All these sheaves are locally free, and the displayed sequence commutes with base change.

\paragraph{The trace convention.}
Take a basis $w_1,\ldots,w_g$ of $W_A$ and a basis $v_1,\ldots,v_g$ of $W_{A^*}$. Let $q_j$ be any lifts to $H_A$ of the basis dual to the $v_j$. We use the order
\begin{equation}\label{eq:hodge-first}
 w_1,\ldots,w_g,q_1,\ldots,q_g.
\end{equation}
The top trace of the exterior product gives a unit
\[
 b=\Tr_A(w_1\wedge\cdots\wedge w_g\wedge q_1\wedge\cdots\wedge q_g).
\]
The result is independent of the lifts of the $q_j$. Define
\begin{equation}\label{eq:beta-def}
 \beta_A:\omega_A\xrightarrow{\sim}\omega_{A^*},\qquad
 w_1\wedge\cdots\wedge w_g\longmapsto b\,v_1\wedge\cdots\wedge v_g.
\end{equation}
This is the determinant of \eqref{eq:hodge}, followed by cup product and top trace. The cup--trace map is an isomorphism of lines: base change to every geometric fibre gives the exterior-cohomology and trace isomorphism, and invertibility of a map between lines is detected on fibres. See \citet[Proposition~7.27 and Corollaries~7.29--7.30, pp.~108--109]{EGM} for the field-level Hodge calculation. In particular, \eqref{eq:beta-def} specifies a map of the lines themselves, including its sign.

For a sheaf $M$ on $S$, write $M^{(p)}=F_S^*M$. Verschiebung induces $V^*:W_A\to W_A^{(p)}$. Its determinant is the Hasse section
\begin{equation}\label{eq:hasse}
 \sigma_A=\det(V^*)\in H^0(S,\omega_A^{\otimes(p-1)}).
\end{equation}
Here we identify $F_S^*\omega_A$ with $\omega_A^{\otimes p}$.

\begin{theorem}[Hasse duality with trace normalization]\label{thm:H}
For every $p$, $S$, and $A/S$ as above, the isomorphism \eqref{eq:beta-def} satisfies
\begin{equation}\label{eq:H}
 \beta_A^{\otimes(p-1)}(\sigma_A)=\sigma_{A^*}.
\end{equation}
This equality holds on $S$ itself, including its nilpotents.
\end{theorem}

Thus the trace comparison carries the Hasse invariant to the dual Hasse invariant, rather than merely carrying one nonvanishing locus to the other. Fargues proved Hasse compatibility for truncated Barsotti--Tate groups on the formal bases of his canonical-subgroup theory \citep[Proposition~2 and proof, author manuscript pp.~8--10]{Fargues}. General Hasse duality for $p$-divisible groups over arbitrary characteristic-$p$ schemes is \citet[Theorem~2.1.4, author manuscript pp.~11--12]{Bijakowski}; his proof uses two complementary-rank subbundles and their determinants. The proof below identifies the abelian comparison with precisely \eqref{eq:beta-def}. It is an exposition of that known duality with the normalization required by the second theorem.

\paragraph{The ordinary determinant.}
Suppose now that $A/S$ is ordinary: every geometric fibre has $p$-rank $g$. The maximal \'etale quotients of $A[p^n]$ form an \'etale $p$-divisible group of height $g$. Its Tate module is the rank-$g$ lisse $\ZZ_p$-local system, understood as the compatible family of finite locally constant $\ZZ/p^n$-sheaves,
\[
 L_A=T_p(A)_{\et}=\varprojlim_n A[p^n]_{\et},
\]
and likewise for $A^*$. A polarization $\lambda:A\to A^*$ has locally constant positive square-root degree
\[
 \delta_\lambda=\sqrt{\deg\lambda}.
\]
On a fibre, if $\lambda=\lambda_L$, Riemann--Roch gives $\delta_\lambda=\chi(L)>0$. Initially define a rational map of determinant local systems by
\begin{equation}\label{eq:alpha-rational}
 \alpha_{A,\lambda}
 =\delta_\lambda^{-1}\det T_p(\lambda)_{\et}:
 \det L_A\otimes\QQ_p\longrightarrow\det L_{A^*}\otimes\QQ_p.
\end{equation}
A polarization may be chosen locally for the \'etale topology; the independence in the next theorem makes this local choice harmless.

The Poincar\'e character pairing and $d\log$ give the ordinary identifications
\begin{equation}\label{eq:ordinary-lines}
 \omega_A=\det(L_{A^*}/pL_{A^*})\otimes_{\FF_p}\OO_S,
 \qquad
 \omega_{A^*}=\det(L_A/pL_A)\otimes_{\FF_p}\OO_S.
\end{equation}
They are specified in Section~\ref{sec:dlog}. For a morphism of $\ZZ_p$-local systems, a bar denotes its reduction modulo $p$.

\begin{theorem}[The ordinary integral determinant and its sign]\label{thm:D}
Let $A/S$ be ordinary with the conventions above, and put
\[
 \epsilon_g=(-1)^{g(g+1)/2}.
\]
The rational map \eqref{eq:alpha-rational} is an integral isomorphism
\[
 \alpha_A:\det L_A\xrightarrow{\sim}\det L_{A^*},
\]
independent of the polarization. These maps glue and commute with base change. Positive Poincar\'e biduality gives
\begin{equation}\label{eq:bidual-main}
 \alpha_{A^*}=\alpha_A^{-1}.
\end{equation}
Under \eqref{eq:ordinary-lines}, the trace comparison is
\begin{equation}\label{eq:D}
 \beta_A=\epsilon_g\,\overline{\alpha_A}^{-1},
 \qquad \beta_A^{-1}=\epsilon_g\,\overline{\alpha_A}.
\end{equation}
All assertions hold for $p=2$ and for polarizations with $p\mid\delta_\lambda$. If $B/S$ is ordinary of relative dimension $h$, the direct-sum determinant identifications give
\begin{equation}\label{eq:product-main}
 \alpha_{A\times B}=\alpha_A\otimes\alpha_B,\qquad
 \beta_{A\times B}=(-1)^{gh}(\beta_A\otimes\beta_B).
\end{equation}
Equivalently, $\epsilon_{g+h}=(-1)^{gh}\epsilon_g\epsilon_h$.
\end{theorem}

In words, Grothendieck's normalized ordinary determinant is integral and canonical even when the polarization has inseparable degree. Its dual is its inverse. The comparison with Hodge determinants has one dimension-dependent sign, and that sign changes under the permutation needed to place all Hodge factors first. For two ordinary elliptic curves in odd characteristic, $\epsilon_1^2=1$ whereas $\epsilon_2=-1$. The sign-free product rule is therefore incompatible with \eqref{eq:hodge-first}.

\paragraph{Credit and neighboring results.}
The determinant divided by $\sqrt{\deg\lambda}$ and the two duality questions come from Grothendieck's manuscript. Bijakowski's general-base Hasse theorem supplies the exact known comparison to which Theorem~\ref{thm:H} belongs. The deformation inputs are \citet[Corollary~4.8(i), pp.~180--181]{Illusie} and the Serre--Tate theorem as stated and proved by \citet[Theorem~1.2.1, pp.~143--146]{Katz}. The ordinary canonical lift and lifting of all homomorphisms are stated by \citet[Theorem~1.4 and the following paragraph, p.~115]{DworkOgus}. Riemann--Roch and the Tate-module degree formula are used at \citet[Theorems~9.11 and 12.18, pp.~131--133 and 186--187]{EGM}. The polynomial argument below proves polarization independence, including its unit, rather than only comparing valuations. The cup-product calculation fixes the sign in Grothendieck's ordinary comparison. No priority claim is made for these formulations.

There is a different determinant problem in \citet[Theorem~1.2 and the surrounding introduction, pp.~1--2]{MaillotRossler}: for a symmetric rigidified ample line bundle, they prove twelfth-power triviality, with sharp exponent, of a line formed from $\det f_*L$ and the Hodge line. That coherent determinant bundle is not the \'etale Tate determinant in \eqref{eq:alpha-rational}. The recent OpenAI preprint \citep[Theorem~1.1 and Introduction, pp.~1--3]{OpenAI} states a rational trace comparison for specialized Hodge classes paired with divisor products. Its normalization is rational and concerns good reduction; it is neighboring context, not an input to the integral ordinary comparison here.

\paragraph{Scope and proof mechanism.}
Theorems~\ref{thm:H} and \ref{thm:D} use different descent arguments. For Hasse duality one must prove an identity over nonreduced deformation rings; equality at geometric points would not suffice. A trace-normalized Frobenius matrix and the complementary-minor identity prove the equality over a formally smooth mixed-characteristic deformation base, and completion then gives every characteristic-$p$ base. For the ordinary determinant, a self-dual kernel proves integrality. A monic polynomial identity eliminates a possible sign in polarization independence. The canonical lift permits the Chern-class and trace calculation in characteristic zero before reducing modulo $p$. Finally, Hasse duality makes the remaining comparison a map of \'etale $\FF_p$-lines, for which geometric-point descent is valid even over a nonreduced base. Ordinarity is essential only for this latter Tate-determinant construction.

\section{Trace-normalized Frobenius and Hasse duality}\label{sec:hasse-proof}

We prove the Hasse identity over a deformation ring that remembers nilpotents, and then descend it to the given base. We use the standard identification of the first crystalline cohomology of an abelian scheme with its contravariant Dieudonn\'e crystal. Evaluation is finite locally free, carries Frobenius and Verschiebung with $FV=VF=p$, and gives \eqref{eq:hodge}. The corresponding Hodge and conjugate filtrations, including their duality conventions, are recalled in \citet[\S2.1, pp.~11--12]{Bijakowski}.

\begin{lemma}[The determinant trace over a formal deformation]\label{lem:trace}
Let $A_0$ be an abelian variety of dimension $g$ over a perfect field $k$. Let
\[
 R=W(k)[[t_1,\ldots,t_{g^2}]]
\]
be the universal formal deformation base of $A_0$. Let $\varphi:R\to R$ be Witt Frobenius on $W(k)$ and satisfy $\varphi(t_i)=t_i^p$. In a basis of its first crystalline cohomology whose ordered top cup product has trace one, the linearized Frobenius has determinant $p^g$.
\end{lemma}
\begin{proof}
The formal deformation exists with this base by \citet[Corollary~4.8(i), pp.~180--181]{Illusie} and \citet[Theorem~1.2.1, pp.~143--146]{Katz}. It is enough to use the formal family; its algebraization is not required. The chosen Frobenius lift preserves the ideal $(t_1,\ldots,t_{g^2})$.

Write $M$ for the finite free $R$-module obtained from the first crystal. Cup product followed by the proper crystalline trace is alternating in its degree-one arguments: graded commutativity kills repeated arguments because the target $R$ has no integer torsion. This also handles $p=2$. It therefore gives a map
\begin{equation}\label{eq:det-trace}
 \tau:\det M\longrightarrow R.
\end{equation}
Under crystalline--de Rham evaluation, this is the top de Rham trace. At the closed fibre its reduction is an isomorphism: the field-level Hodge cohomology algebra and Hodge--de Rham degeneration give the top exterior product, and Serre duality gives its trace \citep[Proposition~7.27 and Corollaries~7.29--7.30, pp.~108--109]{EGM}. Both source and target of \eqref{eq:det-trace} are lines. Its scalar is therefore a unit by Nakayama's lemma. This argument requires no assertion that all relative cohomology groups are free.

The trace is horizontal for the Gauss--Manin connection, and the target has the trivial connection. This normalization can also be read from the de Rham class of the zero section: its trace is one, and its local residue representative is horizontal. Cup product and crystalline Frobenius are horizontal. Thus, if $d\in R$ is the scalar of $\det F$ in \eqref{eq:det-trace}, then its relative differential over $W(k)$ is zero.

The Hodge submodule of $M$ has rank $g$. Its Frobenius image lies in $pM$, since Frobenius on de Rham cohomology kills the Hodge subspace modulo $p$. Consequently $d=p^g a$ for an element $a\in R$. Since $R$ and its module of continuous differentials are $p$-torsion-free, $da=0$. If
\[
 a=\sum_{I\in\ZZ_{\ge0}^{g^2}} a_I t^I,
\]
then differentiation in $t_j$ gives $I_j a_I=0$ for every $I_j>0$. The Witt ring is a characteristic-zero domain, so all coefficients with $I\ne0$ vanish. Hence $a\in W(k)$.

On the fibre $t_1=\cdots=t_{g^2}=0$, the trace identity is
\[
 \tau\circ\bigwedge^{2g}F=p^g\varphi\circ\tau.
\]
For an algebraically closed $k$ this is stated explicitly by \citet[\S6, immediately before Definition~6.1, p.~59]{Ogus1979}; it holds over a perfect $k$ by faithfully flat extension to its algebraic closure. The evaluation of the crystal on any formal lift has this same identity. Thus the constant $a$ is one. This proves the lemma.
\end{proof}

\begin{lemma}[Complementary principal minors]\label{lem:minor}
If $C$ is an invertible $2g\times2g$ matrix over a commutative ring and $C_{22}$ is its lower-right $g\times g$ block, then
\begin{equation}\label{eq:jacobi}
 \det((C^{-1})_{11})=\frac{\det C_{22}}{\det C}.
\end{equation}
There is no sign in this identity.
\end{lemma}
\begin{proof}
For a generic matrix over the fraction field of $\ZZ[c_{ij}]$, its lower-right block is invertible. Block elimination gives
\[
 \det C=\det C_{22}\det(C_{11}-C_{12}C_{22}^{-1}C_{21}),
\]
and the upper-left block of the inverse is the inverse of the Schur complement in parentheses. This proves \eqref{eq:jacobi} there. Replace $C^{-1}$ by $\operatorname{adj}(C)/\det C$ and clear powers of $\det C$. The resulting polynomial identity holds in the integral polynomial ring and hence after every specialization with $\det C$ invertible. The selected row and column subsets are the same, so no nonprincipal-minor permutation occurs.
\end{proof}

\begin{proof}[Proof of Theorem~\ref{thm:H}]
Begin on the formal base of Lemma~\ref{lem:trace}. Choose a basis of $M$ adapted to the lifted Hodge submodule and with top trace one. Use its Frobenius pullback as the source basis. The first $g$ columns of $F$ are divisible by $p$, so
\begin{equation}\label{eq:F-matrix}
 F=C\begin{pmatrix}pI_g&0\\0&I_g\end{pmatrix}
\end{equation}
with $C$ integral. Lemma~\ref{lem:trace} and cancellation of $p^g$ in $R$ give $\det C=1$. In particular $C$ is invertible, and
\begin{equation}\label{eq:V-matrix}
 V=\begin{pmatrix}I_g&0\\0&pI_g\end{pmatrix}C^{-1}.
\end{equation}
Modulo $p$, the Hasse map on the Hodge submodule has determinant $\det((C^{-1})_{11})$. The dual Hasse map is the transpose of Frobenius on the Hodge quotient, whose matrix is $C_{22}$; see \citet[Lemma~2.1.2, p.~11]{Bijakowski}. Transposition does not change the determinant. The chosen volume makes the determinant basis on the dual Hodge line the image of the first basis under $\beta_A$. Lemma~\ref{lem:minor} now gives \eqref{eq:H} on $R/pR$. This also explains why the determinant comparison in Bijakowski's Theorem~2.1.4 agrees here with the trace normalization, after taking the $(p-1)$st tensor power.

For completeness, this formal calculation implies the stated equality for an arbitrary $S$. Work on an affine open and descend the abelian scheme, its dual, the relevant line bundles, and their sections to a finitely generated $\FF_p$-algebra $R_0$. At a closed point of $\Spec R_0$, the residue field $k$ is finite and hence perfect. Its completed local ring has a coefficient field. The restriction of the family to that complete ring is induced by a map from the universal formal deformation ring modulo $p$. Therefore the difference of the two sections in \eqref{eq:H} vanishes after completing at every closed point. Completion of a noetherian local ring is faithfully flat. A nonzero section of an invertible sheaf over a noetherian Jacobson scheme has nonzero support at some closed point. The difference is consequently zero over $R_0$, and base change gives the result over $S$. This descent checks nilpotents as well as reduced points.
\end{proof}

\section{The integral determinant and polarization independence}\label{sec:alpha}

We construct $\alpha_A$ without using the comparison with $\beta_A$, so that the later sign calculation does not enter its definition. Throughout this section $A/S$ is ordinary.

\begin{lemma}[Integrality]\label{lem:integral}
For every polarization $\lambda$, the map $\alpha_{A,\lambda}$ in \eqref{eq:alpha-rational} is an isomorphism of integral rank-one $\ZZ_p$-local systems.
\end{lemma}
\begin{proof}
The assertion can be checked at geometric points, since it is an assertion about lattices in \'etale local systems. Let the point have algebraically closed residue field $k$, and let $K$ be the $p$-primary part of $\ker\lambda$. The polarization pairing makes $K$ Cartier self-dual. Ordinarity splits its connected and \'etale parts; the connected part is of multiplicative type. Cartier duality exchanges these two parts. They therefore have the same order, say $p^a$ each. Thus
\[
 v_p(\deg\lambda)=2a,\qquad v_p(\delta_\lambda)=a.
\]
The map on the free \'etale Tate modules is injective, and its cokernel is the \'etale part of $K$, as a finite $\ZZ_p$-module. Its length is $a$, so
\[
 v_p(\det T_p(\lambda)_{\et})=a.
\]
Dividing the determinant by $\delta_\lambda$ gives a $p$-adic unit. Prime-to-$p$ factors of $\delta_\lambda$ are already units in $\ZZ_p$. Hence the rational determinant map is an integral isomorphism, as asserted.
\end{proof}

\begin{lemma}[Degree on the two ordinary factors]\label{lem:degree}
Let $A/k$ be ordinary over an algebraically closed field of characteristic $p$, let $\mu$ be a polarization, and let $u\in\End^0(A)=\End(A)\otimes\QQ$ be self-adjoint for its Rosati involution:
\[
 u^\vee\mu=\mu u.
\]
Then
\begin{equation}\label{eq:degree-square}
 \det(u\mid L_A\otimes\QQ_p)^2=\deg u.
\end{equation}
For a nonisogeny both sides are zero; for a rational endomorphism degree is extended homogeneously.
\end{lemma}
\begin{proof}
Take the canonical ordinary lift $\widetilde A/W(k)$. Every endomorphism and every homomorphism between ordinary abelian varieties lifts uniquely, and the canonical lift commutes with duality; these assertions hold at $p=2$ as well \citep[Theorem~1.4 and following paragraph, p.~115]{DworkOgus}. Clear a denominator of $u$ if needed. The split connected--\'etale sequence of the canonical lift gives, on its geometric generic fibre, the exact decomposition of rational Tate modules
\[
 V_p(\widetilde A)=
 (L_{A^*}\otimes\QQ_p)^{\vee}(1)\ \oplus\ (L_A\otimes\QQ_p).
\]
The connected summand is the Tate module of the multiplicative part. Functoriality of Cartier duality means that its endomorphism is dual to $u^\vee$ on $L_{A^*}$, not the inverse of that endomorphism. Thus the full determinant is
\[
 \det(u^\vee\mid L_{A^*}\otimes\QQ_p)\,
 \det(u\mid L_A\otimes\QQ_p).
\]
Rosati self-adjointness conjugates the two endomorphisms by the rational isomorphism $T_p(\mu)_{\et}$, so the two determinants are equal.

The generic fibre has characteristic zero. We may therefore apply the ordinary \'etale Tate-module degree theorem with coefficient prime $p$: \citet[Theorem~12.18 and proof, pp.~186--187]{EGM} identifies its full Tate determinant with degree. An integral isogeny that lifts between these abelian schemes has finite flat kernel, and its degree is unchanged by specialization. Clearing denominators gives the same conclusion for a rational isogeny. Nonisogenies follow from the degree polynomial, or by applying the identity to $u+t$ for infinitely many integers $t$ and specializing its polynomial identity at zero; see \citet[Proposition~12.15 and Definition~12.16, p.~185]{EGM}. This proves \eqref{eq:degree-square}. The application of the Tate-module theorem has been made only in characteristic zero, where its prime-to-characteristic hypothesis is satisfied.
\end{proof}

\begin{proposition}[Polarization independence]\label{prop:independence}
The integral map $\alpha_{A,\lambda}$ is independent of $\lambda$. Hence local choices of polarizations define a canonical map $\alpha_A$ commuting with base change.
\end{proposition}
\begin{proof}
Again it suffices to prove equality at a geometric point. Let $\lambda=\lambda_L$ and $\mu=\lambda_M$ be polarizations over its algebraically closed field, and put
\[
 \eta=\mu^{-1}\lambda\in\End^0(A).
\]
Symmetry of $\lambda$ and $\mu$ under positive Poincar\'e biduality gives $\eta^\vee\mu=\mu\eta$. The same is true of $\eta+t$. Define two polynomials over $\QQ_p$ by
\[
 P(t)=\det(\eta_{\et}+tI_g),\qquad
 Q(t)=\frac{\chi(L\otimes M^{\otimes t})}{\chi(M)}.
\]
The second expression denotes its Riemann--Roch polynomial:
\[
 Q(t)=\frac{\deg\bigl((c_1(L)+t c_1(M))^g\bigr)}{g!\,\chi(M)}.
\]
Here $\deg$ applies to the top intersection class. $P$ is monic of degree $g$ by its determinant definition; $Q$ is monic of degree $g$ by \citet[Theorem~9.11, pp.~131--133]{EGM}.

For every integer $t\ge0$, $L\otimes M^{\otimes t}$ is ample, and its polarization is $\lambda+t\mu=\mu(\eta+t)$. Multiplicativity of degree, Lemma~\ref{lem:degree}, and Riemann--Roch give
\begin{align*}
 P(t)^2
 &=\deg(\eta+t)
 =\frac{\deg(\lambda+t\mu)}{\deg\mu}\\
 &=\frac{\chi(L\otimes M^{\otimes t})^2}{\chi(M)^2}
 =Q(t)^2.
\end{align*}
This holds at infinitely many integers, so $P^2=Q^2$ in $\QQ_p[t]$. As this ring is an integral domain, $(P-Q)(P+Q)=0$ implies $P=Q$ or $P=-Q$. Their monic leading coefficients force $P=Q$, also for $p=2$ since the calculation is in $\QQ_2$.

Evaluate at $t=0$. Since $\chi(L)=\delta_\lambda$ and $\chi(M)=\delta_\mu$, we obtain
\[
 \frac{\det\lambda_{\et}}{\det\mu_{\et}}
 =\det\eta_{\et}
 =\frac{\delta_\lambda}{\delta_\mu}.
\]
Consequently $\delta_\lambda^{-1}\det\lambda_{\et}
=\delta_\mu^{-1}\det\mu_{\et}$. This is an equality of the units themselves, not just their valuations.

A polarization exists \'etale locally on the base. The equalities just proved are equalities of morphisms of \'etale local systems, which are detected on geometric stalks. Hence the local integral maps agree on overlaps and glue. The definition on any local polarization commutes with base change, so the glued map does too.
\end{proof}

\begin{proposition}[Positive biduality]\label{prop:bidual}
Under $\iota_A:A\simeq A^{**}$ one has $\alpha_{A^*}=\alpha_A^{-1}$.
\end{proposition}
\begin{proof}
Work locally with a polarization $\lambda$. Choose a positive integer $c$ such that
\[
 \nu=c\lambda^{-1}:A^*\longrightarrow A
\]
is an integral polarization. Here is a verification of existence. An integer killing $\ker\lambda$ makes $c\lambda^{-1}$ integral and symmetric. On a geometric fibre, a symmetric homomorphism is induced by a line bundle $N$ \citep[Proposition~11.2 and Corollary~11.3, pp.~159--160]{EGM}. The pullback $\lambda^*N$ has polarization $c\lambda$, so it is ample; ampleness depends only on the associated polarization \citep[Proposition~11.4, p.~160]{EGM}. Finite surjective pullback detects ampleness, hence $N$ is ample. Thus the integral symmetric isogeny is a polarization on each fibre. Identify $A$ with $A^{**}$ by positive biduality. Then
\[
 \delta_\nu=\frac{c^g}{\delta_\lambda},\qquad
 \det\nu_{\et}=\frac{c^g}{\det\lambda_{\et}}.
\]
Proposition~\ref{prop:independence} permits computing $\alpha_{A^*}$ with $\nu$. The quotient of the two displayed scalars is $\delta_\lambda/\det\lambda_{\et}=\alpha_A^{-1}$. The equality glues. Neither $c$ nor $\delta_\lambda$ needs to be prime to $p$.
\end{proof}

\section{Ordinary differentials and the Chern-class sign}\label{sec:dlog}

We now identify which integral determinant is measured by the top trace, including the sign coming from the de Rham cup product. Let $A/k$ be ordinary over an algebraically closed field, and first work on its canonical lift to $W=W(k)$.

\paragraph{The character normalization.}
The multiplicative part of the ordinary $p$-divisible group of $A$ is the Cartier dual of the \'etale part of $A^*$. A Tate vector $a^*\in L_{A^*}$ determines a character of the formal multiplicative group, and its logarithmic differential is an invariant one-form
\[
 a^*\longmapsto d\log e_A(-,a^*).
\]
This is the ordinary identification in \citet[\S3.3, pp.~161--162]{Katz}. It gives
\begin{equation}\label{eq:dlog}
 W_A\simeq L_{A^*}\otimes_{\ZZ_p}W,
 \qquad W_{A^*}\simeq L_A\otimes_{\ZZ_p}W.
\end{equation}
On a general ordinary characteristic-$p$ base the same character construction is \'etale local and reduces to \eqref{eq:ordinary-lines}. The characters are defined on the multiplicative subgroup; a splitting of the whole ordinary $p$-divisible group over that base is not needed.

For a homomorphism $f:A\to B$, naturality of the Poincar\'e pairing gives
\begin{equation}\label{eq:pair-functor}
 e_B(fx,y)=e_A(x,f^\vee y).
\end{equation}
Taking $d\log$ shows that $f^*$ on differentials is induced by $f^\vee$ on the corresponding dual \'etale Tate module. In particular, because $\lambda^\vee=\lambda$, the matrix of
\[
 \lambda^*:W_{A^*}\longrightarrow W_A
\]
is the matrix of $\lambda_{\et}:L_A\to L_{A^*}$. The inverse produced by reversing the two entries of a Weil pairing must not be inserted again into \eqref{eq:pair-functor}.

Choose bases $e_1,\ldots,e_g$ of $L_A$ and $e_1^*,\ldots,e_g^*$ of $L_{A^*}$. Let $B=(B_{ij})$ be the matrix of $\lambda_{\et}$ in these bases. Let $w_i$ be the forms on $A$ labelled by $e_i^*$, and $v_j$ the forms on $A^*$ labelled by $e_j$. Finally, let $q_j$ be the basis of $\Lie(A^*)=H^1(A,\OO_A)$ dual to the $v_j$, lifted to the unit-root summand of $H^1_{\dR}$. The canonical lift splits the Hodge and unit-root summands, of slopes one and zero respectively.

\begin{lemma}[The mixed Chern class]\label{lem:chern}
Let $L$ induce $\lambda$ on the canonical lift, locally if necessary. In the rational de Rham exterior algebra one has
\begin{equation}\label{eq:chern}
 c_1(L)=-\sum_{i,j}B_{ij}w_i\wedge q_j,
\end{equation}
up to transposing $B$ if the opposite row convention is chosen. In particular the determinant of its mixed coefficient matrix is $(-1)^g\det B$.
\end{lemma}
\begin{proof}
We first identify the coherent Hodge component. Choose a cover trivializing $L$, with transition cocycle $g_{ab}$. For a tangent vector $v$, translation by an infinitesimal point changes this cocycle by
\[
 g_{ab}\longmapsto g_{ab}\bigl(1+\varepsilon\,\iota_v(d\log g_{ab})\bigr).
\]
Consequently the differential of $x\mapsto t_x^*L\otimes L^{-1}$ is represented in $H^1(\OO_A)$ by $\iota_v(d\log g_{ab})$. Thus the tensor in
\[
 H^1(\Omega^1_{A/W})=W_A\otimes H^1(\OO_A)
\]
representing the Hodge component of $c_1(L)$ is exactly the tensor of $d\lambda$. By the functoriality in \eqref{eq:pair-functor}, its coefficient matrix in the chosen bases is $B$, or its transpose. There is no sign in this coherent tensor calculation.

The sign appears when this tensor is expressed as a cup product of degree-one de Rham classes. In the \v Cech--de Rham total complex we use
\[
 D=\delta+(-1)^q d
\]
on \v Cech degree $q$. The product of bidegrees $(a,q)$ and $(a',q')$, with $a$ the form degree, has the factor $(-1)^{a q'}$. A global one-form has bidegree $(1,0)$ and a \v Cech class in $H^1(\OO_A)$ has bidegree $(0,1)$. Their de Rham product is therefore minus their coherent \v Cech product. This gives the minus sign in \eqref{eq:chern} for the mixed component.

There are no other components over $K=\operatorname{Frac}W$. The crystalline first Chern class has Frobenius slope one, since Frobenius pulls a line bundle back to its $p$th tensor power. In the canonical splitting, $\bigwedge^2W_A$ has slope two and $\bigwedge^2$ of the unit-root summand has slope zero. A slope-one class has zero projection to these two summands. The mixed tensor already computed is therefore the entire class. The needed polarization and its line bundle can be chosen on the canonical lift by \citet[Theorem~1.4 and following paragraph, p.~115]{DworkOgus}; making a local choice of the line does not affect its Chern class or the determinant identity.
\end{proof}

\begin{lemma}[The divided exterior power]\label{lem:wedge}
For the ordered basis \eqref{eq:hodge-first} and $B$ as above,
\begin{equation}\label{eq:wedge}
 \frac{c_1(L)^g}{g!}
 =\epsilon_g\det(B)\,
 w_1\wedge\cdots\wedge w_g\wedge q_1\wedge\cdots\wedge q_g
\end{equation}
over $K$, where $\epsilon_g=(-1)^{g(g+1)/2}$.
\end{lemma}
\begin{proof}
Expand the $g$th power of \eqref{eq:chern}. A nonzero term must use each $w_i$ once and each $q_j$ once, hence corresponds to a permutation $\pi\in S_g$. The degree-two factors commute, so the $g!$ orders of these factors have the same sign; this accounts for the denominator. Their coefficients contribute $(-1)^g\prod_iB_{i,\pi(i)}$. Starting from
\[
 (w_1,q_{\pi(1)},w_2,q_{\pi(2)},\ldots,w_g,q_{\pi(g)}),
\]
moving each $q$ past the later $w$'s contributes
$(-1)^{(g-1)+(g-2)+\cdots+1}=(-1)^{g(g-1)/2}$.
Ordering the remaining $q$'s contributes $\operatorname{sgn}(\pi)$.
Summing over $\pi$ gives $\det B$, and the two other signs multiply to $\epsilon_g$. The calculation is in characteristic zero; it does not divide by $g!$ in the residue field.
\end{proof}

\section{The trace scalar, descent, and products}\label{sec:sign}

We finish Theorem~\ref{thm:D} by evaluating the divided Chern class under trace and then explaining why its scalar identity descends through nilpotents.

\begin{proof}[Proof of the comparison in Theorem~\ref{thm:D}]
At an algebraically closed ordinary point, use the canonical lift and bases of Section~\ref{sec:dlog}. Riemann--Roch and the degree formula give
\begin{equation}\label{eq:rr-trace}
 \Tr_A\left(\frac{c_1(L)^g}{g!}\right)
 =\chi(L)=\delta_\lambda;
\end{equation}
see \citet[Theorem~9.11 and proof, pp.~131--133]{EGM}. Let $b$ be the scalar of $\beta_A$ in the two determinant bases. Lemma~\ref{lem:wedge} gives
\[
 \delta_\lambda=\epsilon_g\det(B)b.
\]
We may divide in $K$ and obtain
\begin{equation}\label{eq:scalar}
 b=\epsilon_g\frac{\delta_\lambda}{\det(B)}
   =\epsilon_g\alpha_A^{-1}.
\end{equation}
Both sides are integral units: the left side because the determinant trace is an isomorphism, the right side by Lemma~\ref{lem:integral}. We may therefore reduce the equality modulo $p$, even if $p\mid\delta_\lambda$. The cancellation took place over $K$ before reduction; no zero scalar has been divided out in characteristic $p$.

Now return to an arbitrary ordinary $S$. Under the $d\log$ identifications, the determinant Hasse section is one in every local $\FF_p$-basis: Verschiebung acts as the identity on the character labels, with the indicated Frobenius twist on coefficients. Theorem~\ref{thm:H} consequently says that a local unit $b$ representing $\beta_A$ satisfies
\[
 b^{p-1}=1,\qquad\text{equivalently }b^p=b.
\]
The scheme of such units is the constant finite \'etale group $\FF_p^*$. The morphism $\overline{\alpha_A}^{-1}$ is already a morphism of the same \'etale determinant lines. Thus their ratio, including the constant factor $\epsilon_g$, is a section of a finite \'etale sheaf. Its equality with one can be checked on geometric stalks, where it was proved by \eqref{eq:scalar}. This gives \eqref{eq:D} over all of $S$. This use of geometric points occurs only after the scalar has been shown to be \'etale; it therefore detects the equality on a nonreduced base too. For $p=2$ the group $\FF_2^*$ is trivial, and $\epsilon_g=1$ in the residue field.

Integrality, independence, and biduality were proved in Section~\ref{sec:alpha}. To obtain the product assertion, compute $\alpha_{A\times B}$ using the product polarization. Its degree square root and its \'etale determinant are the products of those for the two factors, so $\alpha_{A\times B}=\alpha_A\otimes\alpha_B$. For $\beta$, the order furnished by the product trace is
\[
 (W_A,Q_A,W_B,Q_B),\qquad Q_A=W_{A^*}^{\vee},\quad Q_B=W_{B^*}^{\vee}.
\]
The Hodge-first order for the product is $(W_A,W_B,Q_A,Q_B)$. Moving the $g$ factors of $Q_A$ past the $h$ factors of $W_B$ contributes $(-1)^{gh}$. Product trace itself is multiplicative, so this proves \eqref{eq:product-main} and completes the theorem.
\end{proof}

\paragraph{An elliptic curve and a product.}
For $g=1$, Lemma~\ref{lem:chern} reads $c_1(L)=-B\,w\wedge q$. Its trace is $\deg L$, so $\beta_A=-\overline{\alpha_A}^{-1}$. With \eqref{eq:polarization-convention}, the degree-one polarization on an elliptic curve is
\[
 x\longmapsto\OO((-x)-(0)),
\]
which is the negative of $x\mapsto\OO((x)-(0))$. This identifies exactly where the dimension-one sign belongs.

For two one-dimensional factors, write $x=w_1\wedge q_1$ and $y=w_2\wedge q_2$. The mixed Chern class with diagonal entries $b_1,b_2$ is $-b_1x-b_2y$. Since $x^2=y^2=0$ and degree-two classes commute,
\[
 \frac{(-b_1x-b_2y)^2}{2}=b_1b_2x\wedge y
 =-b_1b_2\,w_1\wedge w_2\wedge q_1\wedge q_2.
\]
Thus the dimension-two sign is $-1$, while the product of the two dimension-one signs is $1$. Taking a polarization $p\lambda_0$ on an ordinary elliptic curve gives $B=pB_0$ and $\delta=p\delta_0$; the ratio in \eqref{eq:scalar} is unchanged. This is the smallest example of the cancellation needed when the degree is divisible by $p$.

\section{Scope, exact calculations, and method}\label{sec:scope}

We specify the boundary of the ordinary assertion and the role of the finite calculations. Hasse compatibility in Theorem~\ref{thm:H} has no ordinarity hypothesis. The integral construction in Theorem~\ref{thm:D} uses equality of the connected and \'etale orders in a self-dual ordinary kernel, and cannot be extended by simply replacing the rank-$g$ Tate module with the smaller one outside the ordinary locus. For example, for a supersingular elliptic curve and the polarization $p\lambda_0$, the \'etale Tate module has rank zero. Its determinant map is the identity on $\ZZ_p$, but the degree square root is $p$, so the same quotient would be $1/p$, which is not an integral isomorphism. No nonordinary extension of that quotient is asserted. The statements proved here leave no additional hypothesis on the base or on the degree of an ordinary polarization.

The program \emph{Hasse and determinant duality certificates}, available through the \href{\workbenchurl}{workbench}, corroborates the finite algebra in the proofs by two methods. It expands the complete exterior polynomial over the integers in dimensions one through six and compares every monomial with the permutation determinant. It also checks the 64 product shuffles for dimensions one through eight. Its complementary-minor checks use permutation determinants and cofactor inverses over $\FF_p[\varepsilon]/(\varepsilon^2)$: all 144 cases for $p=2,3,5,7$ and $g=1,2,3$ agree, including 106 cases with nonzero nilpotent determinant coefficient. These finite calculations are separate from the proofs for arbitrary dimensions and arbitrary rings given above.

\paragraph{Method and availability.}
The mathematics and text were produced by GPT-6 Astra Pro. The same model checked the full proofs against the cited primary statements, and a second full reading checked the exposition and conventions. Separate exact calculations checked the exterior signs and nilpotent-ring minors. These are same-model checks; no human has checked the mathematics. A separate proof-unseen, statement-only reconstruction of both main statements was completed and fully compared. It independently establishes the all-prime and nonreduced-base assertions, including polarization degrees divisible by $p$. The text and original programs are dedicated to the public domain under CC0~1.0. Editable sources, build instructions, calculations, and detailed check scopes are available in the \href{\workbenchurl}{schemes-and-moduli workbench}; the \href{\recordurl}{mathematics-workbenches index} records versions and their provenance.

\begin{thebibliography}{99}
\bibitem[Bijakowski(2018)]{Bijakowski}
Bijakowski, S. (2018). Duality, refined partial Hasse invariants and the canonical filtration. \emph{Mathematical Research Letters, 25}(4), 1109--1142. \url{https://stephane-bijakowski.perso.math.cnrs.fr/filtration_dual.pdf}. Section and page locators refer to the author manuscript.

\bibitem[Dwork and Ogus(1986)]{DworkOgus}
Dwork, B., \& Ogus, A. (1986). Canonical liftings of Jacobians. \emph{Compositio Mathematica, 58}(1), 111--131. \url{https://www.numdam.org/article/CM_1986__58_1_111_0.pdf}.

\bibitem[Edixhoven et~al.(n.d.)]{EGM}
Edixhoven, B., van der Geer, G., \& Moonen, B. (n.d.). \emph{Abelian varieties} [Book manuscript]. \url{https://van-der-geer.nl/~gerard/AV.pdf}.

\bibitem[Fargues(2011)]{Fargues}
Fargues, L., with the collaboration of Y. Tian. (2011). La filtration canonique des points de torsion des groupes $p$-divisibles. \emph{Annales scientifiques de l'\'Ecole normale sup\'erieure, 44}(6), 905--961. \url{https://doi.org/10.24033/asens.2157}. Author manuscript: \url{https://webusers.imj-prg.fr/~laurent.fargues/canoniqueHN.pdf}.

\bibitem[Grothendieck(n.d.)]{Grothendieck}
Grothendieck, A. (n.d.). \emph{Crystalline-cohomology notes, cote 9} [Manuscript]. Institut Montpelli\'erain Alexander Grothendieck, Universit\'e de Montpellier. \url{https://grothendieck.umontpellier.fr/9.pdf}. Archive sheets 21 and 28; PDF images 22 and 29.

\bibitem[Illusie(1985)]{Illusie}
Illusie, L. (1985). D\'eformations de groupes de Barsotti--Tate. In L. Szpiro (Ed.), \emph{S\'eminaire sur les pinceaux arithm\'etiques: La conjecture de Mordell} (Ast\'erisque 127, pp.~151--198). Soci\'et\'e Math\'ematique de France. \url{https://www.numdam.org/article/AST_1985__127__151_0.pdf}. Author's erratum: \url{https://www.imo.universite-paris-saclay.fr/~luc.illusie/ErrBT.pdf}.

\bibitem[Katz(1981)]{Katz}
Katz, N. M. (1981). Serre--Tate local moduli. In \emph{Surfaces alg\'ebriques} (Lecture Notes in Mathematics 868, pp.~138--202). Springer. \url{https://doi.org/10.1007/BFb0090648}.

\bibitem[Maillot and R\"ossler(2007)]{MaillotRossler}
Maillot, V., \& R\"ossler, D. (2007). \emph{On the determinant bundles of abelian schemes} [Preprint, version 2]. \url{https://arxiv.org/abs/math/0611105}.

\bibitem[Ogus(1979)]{Ogus1979}
Ogus, A. (1979). Supersingular K3 crystals. \emph{Ast\'erisque, 64}, 3--86. \url{https://www.numdam.org/article/AST_1979__64__3_0.pdf}.

\bibitem[OpenAI(2026)]{OpenAI}
OpenAI. (2026, October 7). \emph{Milne's rationality conjecture for abelian varieties} [Preprint]. \url{https://github.com/openai/math/tree/main/preprints/Milnes-rationality-conjecture-for-abelian-varieties-October-7-2026}.
\end{thebibliography}
\end{document}
